One thus obtains a specialization homomorphism on @@M0@@, but with values not in @@M1@@, but in @@M2@@, where @@M3@@ is obtained from @@M4@@ by an extension @@M5@@ of the base fields (which may be assumed to be of finite type when @@M6@@ is universally Japanese). Passing to geometric fibers as explained in 7.16, one likewise obtains a homomorphism
\sgaStableTarget{sga6.e10.eq.7-17-2}@@M7@@
so here again the ``geometric'' specialization homomorphism takes its values not in @@M8@@, but in @@M9@@, where @@M10@@ is obtained from @@M11@@ by the extension of algebraically closed fields @@M12@@.

\sgaStableTarget{sga6.e10.para.7-17-3}\textbf{7.17.3.} In general, the homomorphism @@M13@@ is not an isomorphism (because of the possible presence of ``continuous'' components in the group @@M14@@). It is, however, easy to verify that it is an isomorphism modulo algebraic equivalence (7.11), and also modulo numerical equivalence (4.5.1) when @@M15@@ is smooth and proper. Consequently, using 7.12, one finds that (7.17.2) induces a specialization homomorphism
\sgaStableTarget{sga6.e10.eq.7-17-3-1}@@M16@@
and when @@M17@@ is proper and smooth, provided that specialization homomorphism (7.17.2) is compatible with numerical equivalence (cf. 7.9.4), one likewise obtains a homomorphism
\sgaStableTarget{sga6.e10.eq.7-17-3-2}@@M18@@
We leave to the reader the task of developing (if needed) the corresponding constructions for @@M19@@, @@M20@@, and @@M21@@, and the properties of these specialization homomorphisms that follow trivially from the properties developed in the preceding sections. Let us mention only
